% Cuts of subcritical conformal classes on homotopy four-spheres.
% Initial draft prepared for Edward B. Baker III by Claude (Anthropic, model
% claude-fable-5-2, Cowork session, reasoning effort max as configured),
% 26 November 2026 (America/New_York), at repository state 841d606, from
% papers/master/brainstorm/notes/claude/USES_OF_THE_FIRST_ORDER_AND_MARKED_CUT_TOOLS_20261126.md
% and its checks. Requested by the author when this investigation was set up.
% Revised for Edward B. Baker III on 27 November 2026 by OpenAI Codex.
% Parent replay SP4-CUT-E02 used the draft preserved in 72068ce.
% Parent serving model/effort unavailable; bounded GPT-6.2 Sol/ultra readers.
% Companion results retain their stated hypotheses and review status.
\documentclass[11pt,letterpaper]{article}
\usepackage[margin=1in]{geometry}
\usepackage[T1]{fontenc}
\usepackage{lmodern,microtype,amsmath,amssymb,amsthm,mathtools}
\usepackage[hidelinks]{hyperref}
\setlength{\emergencystretch}{2em}
\numberwithin{equation}{section}
\newtheorem{theorem}{Theorem}[section]
\newtheorem{proposition}[theorem]{Proposition}
\newtheorem{lemma}[theorem]{Lemma}
\newtheorem{corollary}[theorem]{Corollary}
\theoremstyle{definition}
\newtheorem{definition}[theorem]{Definition}
\newtheorem{remark}[theorem]{Remark}
\DeclareMathOperator{\Ric}{Ric}
\DeclareMathOperator{\Rm}{Rm}
\DeclareMathOperator{\tr}{tr}
\DeclareMathOperator{\curl}{curl}
\DeclareMathOperator{\Hess}{Hess}
\newcommand{\cD}{\mathcal D}
\newcommand{\cT}{\mathcal T}
\title{Cuts of subcritical conformal classes\\on homotopy four-spheres}
\author{Edward B. Baker III}
\date{Working manuscript, revised 27 November 2026}

\begin{document}
\maketitle
\begin{abstract}
A positive conformal class on a smooth homotopy four-sphere whose chiral
Weyl energy is below \(8\pi^2\) has a representative with positive scalar
curvature and positive second elementary function of the Schouten tensor.
We study the closed hypersurfaces of such a class. The difference of the
two chiral energies of either side of a hypersurface equals
\(12\pi^2\) times the eta invariant of its conformal structure plus the
integral over it of the trace-free second fundamental form paired with the
magnetic part of the Weyl tensor. Every umbilic hypersurface carries a
metric of positive scalar curvature, and its eta invariant is less than
\(2/3\) in absolute value. For a seam whose two reflected classes have positive Yamabe type,
the number of \(S^1\times S^2\) prime factors is at most one. At free rank
one the two sides have the rational homology of a sphere of dimension
two and a circle, respectively, and the first spends at least \(16\pi^2/3\)
in mean chiral energy. We give a bending criterion for a hypersurface to
supply a positive-scalar-curvature source mapping with nonzero degree to
a specified handle cut. A certified circle-symmetric cap puts a knotted
two-sphere inside an entering example on standard \(S^4\). For a Morse
sweepout a common lower bound on the scalar curvature of all regular
levels excludes index two. Near the round metric such a sweepout can be
chosen through an actual first-width hypersurface. A smooth attained minimum in the
entering band is conformally round if it admits an umbilic round-conformal
three-sphere cut whose two hemisphere-capped classes have positive Yamabe
sign. Producing such a cut or a positive source on an arbitrary marked
entering type remains open.
\end{abstract}

\section{Introduction}

Let \(M\) be a closed, connected, oriented smooth manifold homotopy
equivalent to \(S^4\). Freedman's theorem identifies its homeomorphism
type \cite{Freedman1982}. We ask for a diffeomorphism. In Riemannian
dimension four the Hodge star splits real two-forms into two rank-three
bundles, and the Weyl curvature acts on them by trace-free symmetric
endomorphisms \(W^\pm\) \cite{AHS1978}. For a measurable set \(U\) and a
metric \(g\) put
\[
 E_\pm(g;U)=\int_U|W^\pm_g|_{\rm HS}^2\,dV_g,\qquad
 e(g)=E_+(g;M)=E_-(g;M),\qquad T=4\pi^2 .
\]
The closed energies agree because their difference is \(12\pi^2\tau(M)=0\).
Each chiral measure is conformally invariant. On a region with boundary
the two integrals differ, and that difference is the first subject of
this paper.

The program of the companion papers \cite{Baker2026ChiralMaster} rests on
three facts. A class of positive Yamabe type with \(e<T\) lives on the
standard sphere, by the theorem of Chang, Gursky and Yang
\cite{ChangGurskyYang2003}. A positive class with \(e<2T\) has a
representative with \(R>0\) and \(\sigma_2>0\) for the Schouten tensor
\cite{ChangGurskyYang2002PositiveRicci}; we call such a class
\emph{entering} and such a representative an \emph{entering gauge}. On a
smooth type that carries an entering class the infimum of \(e\) over
positive classes is attained by a smooth Bach-flat class, according to a
working theorem \cite{Baker2026BourAttainment}. The conjecture then
reduces to two statements \cite{Baker2026ChiralMaster}: every smooth
homotopy four-sphere carries an entering class, and

\medskip
\noindent\textup{(S)}\quad
\emph{a smooth homotopy four-sphere with a metric satisfying \(R>0\) and
\(|Z|^2<R^2/12\) is diffeomorphic to \(S^4\).}
\medskip

\noindent
Here \(Z\) is the trace-free Ricci tensor. Neither statement is known.

The constructions tried so far for the first statement assemble a class
from pieces along hypersurfaces, and the arguments proposed for the second
look for a hypersurface along which the manifold can be recognized. Both
uses go back to the collars of \cite{Baker2026CollarProgram} and to the
marked handle cuts of \cite{Baker2026HandleCuts}. This paper studies the
hypersurface itself. We call a closed connected embedded hypersurface a
\emph{cut}. In a homotopy sphere a cut \(\Sigma\) is two-sided and
separates, \(M=V\cup_\Sigma V'\). The conformal class induces on
\(\Sigma\) a conformal metric \([h]\) and a trace-free second fundamental
form \(B^\circ\); the cut is \emph{umbilic} when \(B^\circ=0\).
Our results are the following.

\emph{The ledger} (Section~\ref{sec:ledger}). For a compact oriented
\(V\) with boundary \(\Sigma\),
\[
 E_+(V)-E_-(V)=12\pi^2\bigl(\tau(V)+\eta(\Sigma,[h])\bigr)
 +\oint_\Sigma\langle B^\circ,\curl B^\circ\rangle\,d\sigma ,
\]
and the last integrand is the pairing of \(B^\circ\) with the
magnetic part of \(W\). Hence \(e\ge12\pi^2|\eta(\Sigma,[h])|\) for an
umbilic cut of a homotopy sphere, a bound on the full energy of a collar
by one chiral energy and data on its two ends, and an exact value for
classes that are half conformally flat on each side of an umbilic cut.

\emph{Positivity} (Section~\ref{sec:positivity}). In an entering gauge
every hypersurface satisfies \(R_\Sigma=P+\tfrac23H^2-|B^\circ|^2\) with
\(P>0\). An umbilic cut therefore has positive scalar curvature. This
restricts its topology, turns the recognition theorem of
\cite{Baker2026HandleCuts} into a statement about entering classes, and
excludes conformal reflections whose fixed set carries no metric of
positive scalar curvature.

\emph{Seams and tubes} (Sections~\ref{sec:seams} and~\ref{sec:tube}).
The doubling inequality of \cite{Baker2026ChiralMaster} and a homology
count give the list of levels \(\#_kS^1\times S^2\) that can be umbilic
with positive doubles in an entering class: \(k\le1\). The rank-one case gives a rational homology constraint.
The one-tube description requires an identified normal disk bundle.
For the circle-symmetric construction the closing is standard, including
the certified knotted example in Section~\ref{sec:tube}.

\emph{Sweepouts} (Section~\ref{sec:sweepouts}). A strictly negative
omitted elementary symmetric function of the Hessian forces unbounded
negative scalar curvature near a critical point. Strict positivity of
all four forces divergence to \(+\infty\). Equality can depend on the
third jet. Every index-two point fails the necessary test. This gives
recognition from a uniform lower bound on the regular levels. The
topology of all regular levels of a presentation sphere is compatible
with the same list of three-manifolds even in the presence of index two.

\emph{Free cuts} (Section~\ref{sec:free}). A free spherical cut splits
the half-sum exactly between the two pieces. First-jet-matched round caps
give these energies in the smooth limit. With both capped signs positive,
a smooth minimizer on a hypothetical nonstandard entered type cannot
have such a cut.

The proof reading of Sections~\ref{sec:setting}--\ref{sec:seams},
\ref{sec:sweepouts} and~\ref{sec:free} and the supporting computations
are recorded in \cite{Baker2026MarkedCutCalculations}.
The companion attainment and marked-handle theorems remain working inputs.

\section{Setting}\label{sec:setting}

\subsection{Curvature}
The curvature operator of the unit sphere is the identity. For a metric
\(g\) in dimension four let \(Z=\Ric-\tfrac R4g\) and let
\(A=\tfrac12(\Ric-\tfrac R6g)\) be the Schouten tensor, so that
\[
 \sigma_2(A)=\tfrac12\bigl((\tr A)^2-|A|^2\bigr)
 =\tfrac1{96}R^2-\tfrac18|Z|^2 .
\]
On a closed oriented four-manifold
\begin{equation}\label{eq:cgb}
 8\pi^2\chi=E_++E_-+4\int\sigma_2(A)\,dV,\qquad
 12\pi^2\tau=E_+-E_- .
\end{equation}
The first Pontryagin form of the Levi-Civita connection is
\(p_1=(4\pi^2)^{-1}(|W^+|^2-|W^-|^2)\,dV\).

\begin{lemma}[Entering gauge]\label{lem:gauge}
Let \([g]\) be a class of positive Yamabe type on a homotopy four-sphere
with \(e<2T\). Then \([g]\) contains a metric with \(R>0\) and
\(|Z|^2<R^2/12\). For such a metric and every unit vector \(\nu\),
\[
 0<\Ric(\nu,\nu)<\tfrac R2 .
\]
\end{lemma}
\begin{proof}
By \eqref{eq:cgb} with \(\chi=2\), \(\int\sigma_2(A)=4\pi^2-\tfrac e2>0\).
The theorem of Chang, Gursky and Yang
\cite{ChangGurskyYang2002PositiveRicci} gives a conformal metric with
\(\sigma_2(A)>0\) and \(R>0\) at every point. A trace-free symmetric form
on a four-dimensional Euclidean space has eigenvalues \(z_i\) with
\(\sum z_i=0\), so \(z_1^2\le3\sum_{i>1}z_i^2\) and
\(|z_1|\le\tfrac{\sqrt3}{2}|Z|\). Hence
\(|Z(\nu,\nu)|\le\tfrac{\sqrt3}{2}|Z|<\tfrac R4\), and
\(\Ric(\nu,\nu)=\tfrac R4+Z(\nu,\nu)\) lies strictly between \(0\) and
\(\tfrac R2\).
\end{proof}

We write \(P(\nu)=R-2\Ric(\nu,\nu)\). In an entering gauge \(P>0\) on the
unit sphere bundle. In terms of the Einstein tensor \(G\), \(P=-2G(\nu,\nu)\).

\subsection{Cuts}
Let \(\Sigma\) be a hypersurface with unit normal \(\nu\), induced metric
\(h\), second fundamental form \(B(X,Y)=g(\nabla_X\nu,Y)\), mean curvature
\(H=\tr_hB\) and \(B^\circ=B-\tfrac H3h\). Under \(\hat g=e^{2u}g\),
\begin{equation}\label{eq:conf-B}
 \hat B=e^{u}\bigl(B+(\partial_\nu u)\,h\bigr),\qquad
 \hat B^\circ=e^{u}B^\circ .
\end{equation}
Thus umbilicity is a property of the class, and an umbilic hypersurface is
totally geodesic for a suitable representative. The \emph{Cauchy datum} \(([h],B^\circ)\) uses a covariant
tensor of conformal weight one as in \eqref{eq:conf-B} and a chosen
coorientation. Opposite normal choices negate it. A \emph{free cut}
is an umbilic cut diffeomorphic to \(S^3\) whose conformal structure is
the round one. For an umbilic cut choose a totally geodesic representative on each
side. Reflecting one side \(V\) against a second copy of itself gives
\(\cD V=V\cup_\Sigma\overline V\). In Gaussian coordinates its metric
is \(dt^2+h(|t|)\) with \(h'(0)=0\), of class \(C^{2,1}\) for smooth halves.
We call the cut a \emph{seam with positive doubles} if the reflected
classes on both \(\cD V\) and \(\cD V'\) have strictly positive Yamabe
constant. For any region put \(e_V=(E_+(V)+E_-(V))/2\).
Gluing different smooth pieces with matching first jets generally gives
a \(C^{1,1}\) metric.

Along a hypersurface the Weyl tensor has an electric and a magnetic part,
the trace-free symmetric forms on \(T\Sigma\)
\[
 \mathcal E(X,Y)=W(X,\nu,Y,\nu),\qquad
 \mathcal H_{ij}=\tfrac12\varepsilon_i{}^{kl}W_{klj\nu} .
\]

\section{The chirality ledger}\label{sec:ledger}

For a symmetric form \(B\) on an oriented Riemannian three-manifold let
\((\curl B)_{il}=\varepsilon_l{}^{jk}\nabla_jB_{ik}\). In a local oriented
orthonormal coframe \(e^i\) with \(B_i=B_{ij}e^j\),
\begin{equation}\label{eq:curl-form}
 \langle B,\curl B\rangle\,d\sigma=\sum_iB_i\wedge(d^\nabla B)_i ,
\end{equation}
where \(d^\nabla\) is the covariant exterior derivative on
\(T^*\Sigma\)-valued forms.

\begin{theorem}[Ledger]\label{thm:ledger}
Let \((V,g)\) be a compact oriented Riemannian four-manifold with
boundary \(\Sigma\), oriented by the outward normal followed by an oriented
frame of \(\Sigma\). Let \(\eta(\Sigma,h)\) be the eta invariant of the
odd signature operator, normalized so that the signature theorem of
Atiyah, Patodi and Singer \cite[Theorem 4.14]{APS1975SpectralI} reads
\(\tau(V)=\tfrac13\int_Vp_1-\eta(\Sigma,h)\) for metrics that are products
near the boundary. Then
\begin{equation}\label{eq:ledger}
 E_+(V)-E_-(V)=12\pi^2\bigl(\tau(V)+\eta(\Sigma,h)\bigr)+\cT(\Sigma),
 \qquad
 \cT(\Sigma)=\oint_\Sigma\langle B^\circ,\curl B^\circ\rangle\,d\sigma .
\end{equation}
The term \(\cT(\Sigma)\) depends only on the Cauchy datum
\(([h],B^\circ)\), and \(\eta(\Sigma,h)\) depends only on \([h]\).
\end{theorem}
\begin{proof}
Use the \(g\)-normal collar to choose a product comparison metric
\(g_0\) with the same induced metric and the same unit normal along
\(\Sigma\). Let \(\omega\) and \(\omega_0\) be the
Levi-Civita connections of \(g\) and \(g_0\) on \(TV\), and
\(\theta=\omega-\omega_0\). With \(\Omega_t\) the curvature of
\(\omega_0+t\theta\), \(\tfrac{d}{dt}\tr\Omega_t^2=2\,d\tr(\theta\wedge\Omega_t)\)
and \(\Omega_t=\Omega_0+t\,d_0\theta+t^2\theta\wedge\theta\), so
\[
 p_1(g)-p_1(g_0)=d\,Tp_1,\qquad
 Tp_1=-\frac1{8\pi^2}\Bigl(2\tr(\theta\wedge\Omega_0)
 +\tr(\theta\wedge d_0\theta)+\tfrac23\tr\theta^3\Bigr).
\]
Restrict to \(\Sigma\) and use a frame \((e_1,e_2,e_3,\nu)\). Along
\(\Sigma\) the connection \(\omega_0\) is the Levi-Civita connection of
\(h\) on \(T\Sigma\) and is trivial on the normal line, so \(\Omega_0\)
has only tangential entries, while
\(\theta(X)e_i=-B(X,e_i)\nu\) and \(\theta(X)\nu=B(X,e_i)e_i\) for
\(X\) tangent to \(\Sigma\). Thus \(\theta\) is off-diagonal with
\(\theta^\nu{}_i=-B_i\), \(\theta^i{}_\nu=B_i\). The traces
\(\tr(\theta\wedge\Omega_0)\) and \(\tr\theta^3\) vanish, and
\(\tr(\theta\wedge d_0\theta)=-2\sum_iB_i\wedge(d^\nabla B)_i\). By
Stokes' theorem and \eqref{eq:curl-form},
\[
 \int_Vp_1(g)=\int_Vp_1(g_0)+\frac1{4\pi^2}\oint_\Sigma
 \langle B,\curl B\rangle\,d\sigma .
\]
The signature theorem for \(g_0\) and \(p_1=(4\pi^2)^{-1}(|W^+|^2-|W^-|^2)dV\)
give \eqref{eq:ledger} with \(B\) in place of \(B^\circ\). Write
\(B=B^\circ+\tfrac H3h\). The cross terms in
\(\langle B,\curl B\rangle\) are \(\tfrac H3\varepsilon^{ijk}\nabla_jB^\circ_{ik}\)
and \(\tfrac13\varepsilon^{lji}B^\circ_{il}\nabla_jH\), which vanish by
symmetry, and the term quadratic in \(H\) vanishes as well. So
\(\langle B,\curl B\rangle=\langle B^\circ,\curl B^\circ\rangle\).

Under \(\hat g=e^{2u}g\) the one-forms \(B^\circ_i\) taken in the rescaled
frame are unchanged by \eqref{eq:conf-B}, and the change of tangential connection is
\(\widehat\omega^i{}_j-\omega^i{}_j=u_je^i-u_ie^j\) for
\(\nabla e_j=\omega^i{}_j e_i\). Both resulting terms in the contracted
three-form vanish because \(\sum_jB_j^\circ\wedge e^j=0\). So \(\cT\) is
conformally invariant. Finally let \(h_1=e^{2w}h_0\) and apply
\eqref{eq:ledger} to \(\Sigma\times[0,1]\) with a metric conformal to
\(dt^2+h_0\) that induces \(h_0\) and \(h_1\) on the two ends. In the extended product, reflection about each slice fixes it
pointwise and interchanges the chiralities. Equal Weyl norms survive
conformal scaling, the ends are umbilic, and the signature is zero. Hence \(\eta(\Sigma,h_1)=\eta(\Sigma,h_0)\).
\end{proof}

The conformal invariance of \(\eta\) is due to Atiyah, Patodi and Singer
\cite{APS1975SpectralII}. The sign can also be fixed without using
\eqref{eq:ledger} to compute the boundary eta. On the radius-\(r\)
Fubini--Study ball of holomorphic sectional curvature four, complex
orientation and \(0<r<\pi/2\), the bulk difference is
\(12\pi^2\sin^4r\). Intrinsic transgression along a product-ended path of
Berger metrics gives \(12\pi^2\eta=8\pi^2\sin^4r\). The outward shape
operator gives \(\cT=+4\pi^2\sin^4r\)
\cite[LED, Fubini--Study balls]{Baker2026MarkedCutCalculations}.

\begin{proposition}[Pointwise form]\label{prop:EH}
Use \(R_{abcd}=\langle R^\nabla(e_a,e_b)e_d,e_c\rangle\) and
\((\nu,e_1,e_2,e_3)\) positive. With the magnetic tensor defined above,
\[
 W^\pm \longleftrightarrow \mathcal E\mp\mathcal H,\qquad
 |W^+|^2-|W^-|^2=-4\langle\mathcal E,\mathcal H\rangle,\qquad
 \mathcal H=\operatorname{sym}(\curl B^\circ).
\]
Thus \(\cT(\Sigma)=\oint\langle B^\circ,\mathcal H\rangle\,d\sigma\);
at an umbilic hypersurface the two chiral norms agree.
\end{proposition}
\begin{proof}
Set \(u_i=e^i\wedge\nu^\flat\) and
\(v_i=\tfrac12\varepsilon_i{}^{jk}e^j\wedge e^k\). Then \(*u_i=-v_i\).
In the two triples \(u,v\) the diagonal blocks of the Weyl operator are
both \(\mathcal E\), since it commutes with star, and the off-diagonal
blocks are \(\mathcal H\) and its transpose. In particular
\(W_{ijkl}=+\varepsilon_{ijm}\varepsilon_{kln}\mathcal E^{mn}\).
The chiral bases \((u_i\mp v_i)/\sqrt2\) give the stated difference.
Codazzi in this convention reads
\(R_{klj\nu}=\nabla_k B_{lj}-\nabla_l B_{kj}\). The mixed Schouten
term contributes only to the antisymmetric part after contraction with
\(\varepsilon_i{}^{kl}\), as does the curl of the pure trace of \(B\).
\end{proof}

\begin{corollary}[Eta bound]\label{cor:eta}
Let \(\Sigma\) be an umbilic cut of a class on a homotopy four-sphere
\(M=V\cup_\Sigma V'\). Then
\(E_+(V)-E_-(V)=12\pi^2\eta(\Sigma,[h])=-(E_+(V')-E_-(V'))\) and
\[
 e\ \ge\ 12\pi^2\,|\eta(\Sigma,[h])|=3T\,|\eta(\Sigma,[h])| .
\]
In an entering class \(|\eta(\Sigma,[h])|<\tfrac23\) for every umbilic
cut, and \(|\eta|<\tfrac13\) if \(e<T\).
\end{corollary}
\begin{proof}
Compact interior representatives of two classes in \(H_2(V;\mathbb R)\)
have the same algebraic intersection in \(V\) and in \(M\). Since
\(H_2(M)=0\), the form on the image of absolute in relative homology
has signature zero. Thus
\(\tau(V)=\tau(V')=0\), and \(\cT=0\). If \(\eta\ge0\) then
\(E_+(V)\ge12\pi^2\eta\); otherwise \(E_+(V')\ge12\pi^2|\eta|\). In both
cases \(e=E_+(V)+E_+(V')\ge12\pi^2|\eta|\).
\end{proof}

For the quotient of the round sphere by a finite subgroup \(\Gamma\) of
\(SU(2)\) acting by left multiplication, the fixed point formula gives
\(|\eta|=|\Gamma|^{-1}\sum_{\gamma\ne1}\cot^2(\theta_\gamma/2)\) with
\(\cos\theta_\gamma=\operatorname{Re}\gamma\): the values are \(0\),
\(\tfrac29\), \(\tfrac12\), \(\tfrac45\) for the cyclic groups of orders
two to five, \(\tfrac34\) for the quaternion group and \(\tfrac{49}{36}\)
for the binary tetrahedral group. Here the averaging uses the equivariant signature defect
\cite[Proposition 2.12 and its application]{APS1975SpectralII}.
For a smooth path \(h_s\) and
\(Y_{ij}=\varepsilon_i{}^{kl}\nabla_k(\Ric_{lj}-R_hh_{lj}/4)\),
intrinsic product-ended transgression and integration by parts give
\(12\pi^2\dot\eta=\int\langle\dot h,Y\rangle\,d\mu_h\)
\cite[LED, eta variation]{Baker2026MarkedCutCalculations}.
The round quaternionic link, for which
the bound is \(9\pi^2\), is treated by the same formula in
\cite{Baker2026MarkedNecks}; Corollary~\ref{cor:eta} extends that test to
an arbitrary conformal structure on an arbitrary umbilic cut.

More generally let \(D=12\pi^2\eta+\cT\) for any cut, with the
outward orientation of \(V\). For \(D\ge0\),
\[
                 e-|D|=E_-(V)+E_+(V').
\]
For \(D\le0\) the two chiralities exchange roles. Thus a proposed marked
cut in an entering class must satisfy \(|D|<2T\), including when it is
nonumbilic.

\begin{corollary}[Collars]\label{cor:collar}
For a marked oriented product \(C=\Sigma\times[0,1]\), orient both
slices by the increasing normal. Then
\begin{align*}
 E_-(C)-E_+(C)&=-12\pi^2\big(\eta(h_1)-\eta(h_0)\big)
                       -\cT_1+\cT_0,\\
 |\cT_i|&\le\|B_i^\circ\|_{L^3}
       \|\operatorname{sym}(\curl B_i^\circ)\|_{L^{3/2}}.
\end{align*}
\end{corollary}
The product has zero signature. The estimate is H\"older's inequality
using Proposition~\ref{prop:EH}; each factor is conformally invariant
with the indicated exponent. A bound with the full curl also holds in
a chosen gauge. On fixed-width marked end windows \(C^2\)-close to the
same unit round product by \(\delta_i\), one has
\(|E_-(C)-E_+(C)|=O(\delta_0^2+\delta_1^2)\). Indeed \(B_i^\circ\) is
\(C^1\)-small, and an intrinsic product-ended interpolation from the
round model to \(h_i\) has Weyl blocks \(O(\delta_i)\) on a fixed
strip. Its signature formula bounds the change of eta quadratically.
This comparison applies on the buffered smooth collars specified in
\cite{Baker2026MarkedNecks,Baker2026BourAttainment}. On a Dirac strip
one must supply the end estimates in the same marking in addition to the
spinor estimate \cite{Baker2026DiracTransmission}.

\begin{corollary}[Half-flat pieces]\label{cor:halfflat}
If \(W^+=0\) on \(V\) and \(\partial V\) is umbilic, then
\(e_V=6\pi^2|\tau(V)+\eta(\partial V)|\). A class on a homotopy four-sphere
with \(W^-=0\) on one side of an umbilic cut and \(W^+=0\) on the other
has \(e=12\pi^2|\eta(\Sigma,[h])|\).
\end{corollary}

For smooth connected \(M\), equality in \(e\ge |D|\) forces
both complementary chiral blocks to vanish on their respective open
sides. Each half-flat side is Bach-flat by compactly supported first
variation of the Weyl functional and the signature identity. By
smoothness the Bach tensor vanishes on \(M\). In conformal harmonic
coordinates the determinant-normalized metric is analytic
\cite[Theorem 3.7(c)]{LassasLiimatainen2019ConformalHarmonic}.
The Hodge projections are analytic there, so each vanishing block
continues across the cut and over the connected manifold. The developing
map of a closed simply connected conformally flat four-manifold is a
covering of \(S^4\), hence a diffeomorphism
\cite{SchoenYau1988Conformal}. Thus equality on a smooth homotopy
sphere forces the round conformal class. First-jet-matched pieces can
have positive limiting energy at equality.

\begin{remark}[Cauchy data]\label{rem:cauchy}
LeBrun's local theorem gives the real-analytic half-flat Riemannian
germ with specified analytic conformal metric and second conformal
form \cite[Theorem 3.6.1]{LeBrun1980ComplexGeodesics}.
The Einstein construction at conformal infinity
\cite{LeBrun1982HSpace} has different boundary data.
Biquard's near-round fillability criterion concerns the twistor data
of a whole ball \cite{Biquard2002Autoduales}. Matching analytic germs
or separately finding fillings still leaves the global marking, energy
and positivity of an assembled class to be proved.
\end{remark}

\section{Positivity of cuts}\label{sec:positivity}

\begin{theorem}[Positivity]\label{thm:positivity}
Let \(g\) be an entering gauge and \(\Sigma\) a hypersurface with unit
normal \(\nu\). Then
\begin{equation}\label{eq:gauss}
 R_\Sigma=P(\nu)+\tfrac23H^2-|B^\circ|^2,\qquad P(\nu)>0 .
\end{equation}
Consequently:
\begin{enumerate}
\item if \(|B^\circ|^2\le\tfrac23H^2\), that is \(\sigma_2(B)\ge0\), then
the induced metric has positive scalar curvature;
\item every umbilic hypersurface of an entering class carries a metric of
positive scalar curvature;
\item if \(\Sigma\) is closed and
\(\|B^\circ\|_{L^3(\Sigma)}^2<\inf_{u\ne0}\int_\Sigma\bigl(8|\nabla u|^2
+(P+\tfrac23H^2)u^2\bigr)\big/\|u\|_{L^6}^2\), then \(\Sigma\) has positive
Yamabe type.
\end{enumerate}
\end{theorem}
\begin{proof}
The Gauss equation is \(R_\Sigma=R-2\Ric(\nu,\nu)+H^2-|B|^2\), and
\(H^2-|B|^2=\tfrac23H^2-|B^\circ|^2=2\sigma_2(B)\). Lemma~\ref{lem:gauge}
gives \(P>0\). This proves \eqref{eq:gauss} and (1). An umbilic
hypersurface is umbilic for the entering gauge by \eqref{eq:conf-B}, which
gives (2). For (3), the conformal Laplacian \(-8\Delta+R_\Sigma\) of
\(\Sigma\) has quadratic form
\(\int(8|\nabla u|^2+(P+\tfrac23H^2)u^2)-\int|B^\circ|^2u^2\), and
\(\int|B^\circ|^2u^2\le\|B^\circ\|_{L^3}^2\|u\|_{L^6}^2\).
\end{proof}

\begin{remark}\label{rem:two-convex}
If \(\sigma_1(B)>0\) and \(\sigma_2(B)>0\) then the sum of any two
principal curvatures is positive, because
\(\partial\sigma_2/\partial\kappa_i=\sigma_1-\kappa_i\) is positive on the
cone \(\{\sigma_1>0,\sigma_2>0\}\). In dimension three the hypersurfaces of
case (1) with \(H>0\) are weakly two-convex by continuity.
Strict two-convexity follows if \(\sigma_2(B)>0\);
the eigenvalues \((H,0,0)\) realize the distinction.
\end{remark}

\begin{corollary}[Topology of umbilic cuts]\label{cor:topology}
A closed umbilic hypersurface of an entering class is a connected sum of
spherical space forms and copies of \(S^1\times S^2\). If it is the
boundary of a tubular neighbourhood of an embedded closed connected surface
\(F\), then \(F\) is a sphere or a projective plane.
\end{corollary}
\begin{proof}
The first statement is the classification of closed orientable
three-manifolds of positive scalar curvature
\cite{GromovLawson1980Surgery,Perelman2003Surgery}. The boundary of a
tubular neighbourhood of \(F\) is a circle bundle with orientable
total space; fibre orientation can reverse over a nonorientable base.
For \(\chi(F)\le0\) its pullback to the contractible universal cover
of \(F\) is a product, so its total space is aspherical. This excludes
the listed positive-scalar-curvature types.
\end{proof}

\begin{corollary}[Recognition at a marked cut]\label{cor:marked}
Fix a reduced mixed presentation with two middle handles, winding row
\((1,0)\), and its original upper closing as in \cite{Baker2026HandleCuts}.
Write \(Z_a\) for its winding-one-first cut and \(Z_b\) for the other order.
If a closed oriented positive-scalar-curvature three-manifold maps with
nonzero degree to \(Z_a\) or \(Z_b\), the corresponding cut is
\(S^3\) or \(\#_2(S^1\times S^2)\), respectively, and \(M\cong S^4\).
In an entering class a source is supplied by an umbilic hypersurface,
or by a closed cut satisfying Theorem~\ref{thm:positivity}(3), provided
the indicated degree map is supplied.
\end{corollary}
\begin{proof}
The geometric hypotheses give a positive metric in the induced conformal
class. The degree condition is precisely PSC domination in the companion
theorem on the two normalized mixed cuts \cite{Baker2026HandleCuts}.
An oriented identification with the actual cut has degree \(\pm1\).
\end{proof}

\begin{corollary}[Conformal reflections]\label{cor:reflection}
Let \(\rho\) be an orientation-reversing involution of \(M\) whose fixed
set is a hypersurface \(\Sigma\). If an entering class is
\(\rho\)-invariant, then \(\Sigma\) admits positive scalar curvature. In
particular the double of a compact contractible four-manifold whose
boundary admits no metric of positive scalar curvature has no
reflection-invariant entering class.
\end{corollary}
\begin{proof}
Averaging a representative gives a \(\rho\)-invariant metric in the class,
for which the fixed set is totally geodesic. So \(\Sigma\) is umbilic.
\end{proof}

There is a simultaneous numerical threshold for the cuts of any
\emph{one} smooth entering gauge. Fix a smooth isometric embedding
\(F:(M,g)\longrightarrow\mathbb R^{3+m}\), \(m\ge2\)
\cite[Part C, Theorem 2]{Nash1956Embedding}. Write \(p=\min P>0\) over
unit normals, \(K=\max_L|\tr_L\mathrm{II}_F|\) over tangent three-planes,
and \(\omega_j\) for Euclidean unit ball volumes. Then
\begin{equation}\label{eq:bending}
 Y(\Sigma,[h])\ \ge\ C(g,F)-\|B^\circ\|_{L^3(h)}^2,\quad
 C(g,F)=\frac{b_m^2}{\max\{2,K^2/p\}},\qquad
 b_m=3\left(\frac{(3+m)\omega_{3+m}}{m\omega_m}\right)^{1/3}.
\end{equation}
Indeed the Euclidean mean curvature vector of the cut has squared norm
\(H^2+|\tr_{T\Sigma}\mathrm{II}_F|^2\). Apply Brendle's inequality to
\(u^4+\epsilon\), send \(\epsilon\) to zero, and use Cauchy--Schwarz
\cite[Theorem 1]{Brendle2021Isoperimetric}. This gives
\(b_m^2\|u\|_6^2\le\int(16|\nabla u|^2+(H^2+K^2)u^2)\).
Comparison with the numerator of Theorem~\ref{thm:positivity}(3)
proves \eqref{eq:bending}. Bending and the induced Yamabe constant
are conformal invariants. The displayed certificate is obtained in the
chosen gauge and embedding.

The ambient embedding can be eliminated on a specified physical band.
The following constant depends only on measurements in its charts.

\begin{proposition}[A finite-atlas criterion on a native band]\label{prop:atlas-band}
Let \(h_t\) be the metrics on the three-dimensional slices of a smooth
oriented product band \(Z\times[a,b]\). In finitely many coordinate
charts let \(mI\le h_t\le MI\) on the supports of a quadratic partition
\(\sum_j\chi_j^2=1\), and suppose
\(\sum_j |d\chi_j|_{h_t}^2\le L\). Take a Euclidean Dirichlet constant
\(S_D>0\) such that \(\int|Dv|^2\ge S_D\|v\|_6^2\) on each chart.
In the physical metric put \(Q_t=R-2\Ric(\nu,\nu)+2H^2/3\). For \(q_0>0\) set
\[
 C_{\rm loc}=\frac{S_D\,m/M}{\max\{1/8,L/q_0\}},\qquad
 d_t=\|(q_0-Q_t)_+\|_{L^{3/2}(h_t)},\qquad
 b_t=\int_Z|B_t^\circ|^3\,d\mu_{h_t}.
\]
Then \(Y(Z,[h_t])\ge C_{\rm loc}-d_t-b_t^{2/3}\).
In particular, for any probability measure \(\mu\) on the height interval,
\begin{equation}\label{eq:atlas-band}
 \int b_t\,d\mu(t) <
 \int (C_{\rm loc}-d_t)_+^{3/2}\,d\mu(t)
\end{equation}
implies that some slice has positive Yamabe type. In a native product
marking that slice projects with degree one onto the original cut.
\end{proposition}
\begin{proof}
For the eigenvalues of \(h_t\), ordered increasingly, the gradient density
has least eigenvalue \(\sqrt{\lambda_1\lambda_2/\lambda_3}\ge m/\sqrt M\).
The squared \(L^6\) norm acquires a factor of at most \(\sqrt M\).
The Dirichlet inequality in each chart therefore has constant at least
\(S_D m/M\). Apply it to \(\chi_j u\) and sum. The triangle inequality
applied to \(u^2=\sum_j(\chi_j u)^2\) in \(L^3\), and
\(\sum_j\chi_jd\chi_j=0\), give
\[
 (S_D m/M)\|u\|_6^2
 \le \int_Z \bigl(|du|_{h_t}^2+Lu^2\bigr)\,d\mu_{h_t}.
\]
Compare with \(8|du|^2+q_0u^2\), then pay the two errors by
H\"older in the Gauss identity \(R_{h_t}=Q_t-|B^\circ_t|^2\).
A slice of nonpositive Yamabe sign thus satisfies
\(b_t\ge (C_{\rm loc}-d_t)_+^{3/2}\). Integrating proves the claim.
\end{proof}

If \(C_{\rm loc}-d_t\ge c>0\), the proportion of positive slices is
at least \(1-c^{-3/2}\int b_t\,d\mu\). The native orientation is
increasing-normal first. Relabeling the heights preserves the test
when the measure is transported with the label. For uniform height
measure the calculation is explicit. On a native collar
\(Z\times[a,b]\) let \(h_t\) be the slice metric, \(N>0\) the lapse and
\(X\) the tangential shift, so the unit normal is
\(N^{-1}(\partial_t-X)\). If
\begin{equation}\label{eq:marked-average}
 \frac{1}{8(b-a)}\int_a^b\!\!\int_Z
  N^{-3}|(\partial_t h_t-\mathcal L_X h_t)^\circ|_{h_t}^3
                \,d\mu_{h_t}\,dt < C(g,F)^{3/2},
\end{equation}
some slice has positive Yamabe type and its projection has degree one.
Here \(2NB=\partial_t h_t-\mathcal L_Xh_t\).
The left side equals \(\int b_t\,dt/(b-a)\), so one may also use the
right side of \eqref{eq:atlas-band} from the finite atlas.
When integrating instead against the four-volume of this band the
lapse factor is \(N^{-4}\), since \(dV_g=N\,dt\,d\mu_{h_t}\). More generally a closed oriented
separator homologous in a band to its incoming face maps with total
degree one under any map to \(Z\) restricting to degree one on that face.
At least one component has nonzero degree. Its positive sign must be
verified, or verified for every component. The orientable product band
also admits the intrinsic width test of R\"ade
\cite[Corollary 2.25]{Raede2022Bands}: if \(R\ge\rho>0\) and its intrinsic
end distance \(d\) satisfies \(\rho d^2\ge3\pi^2\), \(Z\) admits positive
scalar curvature. This follows by scaling the four-dimensional band to
scalar curvature at least twelve and using the strict width bound
\(\pi/2\) for a nonpositive-type target.

On the split \(K\)-first diagram of \cite{Baker2026HandleCuts},
the physical latitude metric is \(dq^2+\cos^2 q\,\gamma\),
\(|q|\le\pi/6\), with \(\gamma\) unit round on \(S^3\).
Here \(Q_q=6\sec^2q\) and \(b_q=0\). Two stereographic charts
give \(C_{\rm loc}\ge81/(800\pi^2)\).
There are genuinely bent examples on the same marked diagram. Put
\(z=y_1\) on \(S^3\), let \(\phi_q\) be the time-\(\theta(q)\) flow of
\(z\nabla_\gamma z\), and take
\(h_q=\cos^2 q\,\phi_q^*\gamma\) with physical lapse one and shift zero.
Use a nonconstant \(\theta\) vanishing near both end faces, with
\(|\theta|\le(\log2)/2\). Since
\(\mathcal L_{z\nabla z}\gamma=2\,dz\otimes dz-2z^2\gamma\),
\[
 |B^\circ|=\sqrt{2/3}\,|\theta'|(1-z^2)\circ\phi_q,\qquad
 Q_q=6\sec^2 q+|B^\circ|^2 .
\]
The two stereographic charts on the unperturbed latitudes have
\(m=3/25,\ M=4,\ L=4\pi^2/3\). The flow changes the metric by a
factor between \(1/2\) and \(2\), so \(m=3/50,\ M=8,\ L=8\pi^2/3\)
work for the bent slices. Taking \(q_0=6,\ S_D=3/4\) gives
\(C_{\rm loc}=81/(6400\pi^2)\). From
\(\int_{S^3}(1-z^2)^3\,d\mu_\gamma=35\pi^2/32\) one gets the exact cost
\[
 b_q=\frac{35\pi^2}{32}(2/3)^{3/2}\cos^3 q\,|\theta'(q)|^3.
\]
A nonzero bump with sufficiently small amplitude thus satisfies
\eqref{eq:atlas-band} with strict margin on the original marked band.
This fits the original smooth closing, and small multiples of a fixed
bump retain a closed entering class by \(C^2\) convergence to round.
A change of angular coordinates on the round metric would also change
the shift and would give \(B^\circ=0\)
\cite[finite-atlas estimate and split calibrations]{Baker2026MarkedCutCalculations}.

On the round unit sphere \(p=6\), \(K=3\), \(m=2\), and
\(C^{3/2}=9\sqrt2\,\pi\). For every embedded copy of the winding-one
Mazur boundary \(\Sigma(2,5,7)\) in Isoshima's \(n=1\) double,
\(\int|B^\circ|^3\ge9\sqrt2\,\pi\) in the pulled-back round class.
Its nontrivial fundamental group and smooth contractible filling exclude
positive scalar curvature \cite[Corollary 2.3 and Proposition 4.2]{ChodoshMaximoMukherjee2024Complete};
the marked double and its two orders are identified in
\cite[Figure 37]{Isoshima2025MazurDoubles} and \cite{Baker2026HandleCuts}. In the other order the marked
meridian-first cut is \(\#_2(S^1\times S^2)\). In the same round class
an embedded positive source maps to that target with degree one by the
separate construction in \cite{Baker2026ChiralMaster}.
On the retained equator in the regular and long-cuff minus
Miller--Schupp examples, both complementary groups are \(BS(2,3)\).
Its first Betti number two rules out identification with the winding-one
cut, and each complement of the actual winding-zero cut in a normalized
mixed presentation has group \(\mathbb Z\). Thus this embedded equator
cannot be identified with either of these native cuts. A distinct degree
map or a physical witness on an untreated presentation remains to be
supplied \cite{Baker2026ChiralMaster}.

For the native Gompf--Scharlemann--Thompson links
\(L_{n,1}=Q\cup V_n\), \(n\ge3\), neither single zero-surgery cut can
receive a nonzero-degree map from a PSC three-manifold. Domination
would force \(S^1\times S^2\) by the zero-matrix surgery trace in
\cite{Baker2026HandleCuts}. The fundamental group of \(S^1\times S^2\) could not surject onto
the complementary one-relator quotient, which itself maps onto
\(C_2*C_3\) or \(C_{n+1}*C_n\). The marked word estimates exclude
both labels through three unstabilized reverse slides. At \(n=3\)
the stationary fourth label and \(224\) of \(256\) signed moving-label
histories are excluded by the recorded sample tests
\cite[SLIDE, cut certificates]{Baker2026MarkedCutCalculations}.
There is also a marked exterior on the standard sphere for each native
proper slice disk. Doubling that disk and folding its exterior preserves
the disk exterior group and the recorded meridian. The product filling
recovers \(S^4\). The relative handle calculation also identifies
the twisted filling of each of these two disk doubles with \(S^4\).
A chiral bound would require a metric on the specified exterior
\cite[SLIDE, native disk doubles]{Baker2026MarkedCutCalculations}.

\section{Seams}\label{sec:seams}

We recall the doubling inequality of \cite{Baker2026ChiralMaster}. If
\(\Sigma=\partial V\) is umbilic, the doubled class on \(\cD V\) satisfies
\(E_\pm(\cD V)=E_+(V)+E_-(V)=2e_V\), and \(\tau(\cD V)=0\),
\(\chi(\cD V)=2\chi(V)-\chi(\Sigma)=2\chi(V)\). Gursky's inequalities for
positive classes \cite{Gursky1998WeylFunctional,Sung2021Weyl}, applied to smooth
approximations of the doubled class, give
\begin{equation}\label{eq:doubling}
 e_V\ge T\chi(V)\ \ \text{if } b_1(\cD V)>0,\qquad
 e_V\ge\tfrac23T\chi(V)\ \ \text{if } b_2(\cD V)>0,
\end{equation}
whenever the doubled class has positive Yamabe type.

\begin{lemma}\label{lem:homology}
Let \(\Sigma\) be a connected cut of a homotopy four-sphere, \(k=b_1(\Sigma)\),
\(M=V\cup_\Sigma V'\), and \(a=b_1(V)\). Then \(b_1(V')=k-a\),
\(b_2(V)=k-a\), \(b_3(V)=0\), \(\chi(V)=1+k-2a\), and
\(b_1(\cD V)=a\), \(b_2(\cD V)=2(k-a)\).
\end{lemma}
\begin{proof}
The Mayer--Vietoris sequence of \(M=V\cup V'\) and the vanishing of
\(H_j(M)\) for \(0<j<4\) give \(H_j(\Sigma)\cong H_j(V)\oplus H_j(V')\)
for \(j=1,2\) and \(H_3(V)=0\). So \(H_1(\Sigma)\to H_1(V)\) is onto with
kernel of rank \(k-a\). The kernels of \(H_1(\Sigma)\to H_1(V)\) and
\(H_2(\Sigma)\to H_2(V)\) are annihilators of each other for the
intersection pairing of \(\Sigma\), over the rationals; so the second
kernel has rank \(a\) and \(b_2(V)=k-a\). For the double, the
Mayer--Vietoris map \(H_1(\Sigma)\to H_1(V)\oplus H_1(V)\) has rank \(a\),
which gives \(b_1(\cD V)=a\), and \(\chi(\cD V)=2\chi(V)\) gives \(b_2\).
\end{proof}

\begin{theorem}[Seam atlas]\label{thm:atlas}
Let \(\Sigma\) be a connected seam with positive doubles of a
class on a homotopy four-sphere. Put \(k=b_1(\Sigma)\) and \(a=b_1(V)\).
\begin{enumerate}
\item If \(0<a<k\) then \(e\ge2T\).
\item If \(a\in\{0,k\}\) and \(k\ge1\) then \(e\ge\tfrac23T(1+k)\).
\end{enumerate}
Hence an entering class has such a seam only for \(k\le1\), and for
\(k=1\) the side with \(b_2=1\) has \(e_V\ge\tfrac43T\) and the other side
has \(e_{V'}<\tfrac23T\).
\end{theorem}
\begin{proof}
In case (1), \(b_1(\cD V)=a>0\) and \(b_1(\cD V')=k-a>0\), so by
\eqref{eq:doubling} \(e_V\ge T\chi(V)\) and \(e_{V'}\ge T\chi(V')\), and
\(\chi(V)+\chi(V')=2\). If one of the Euler characteristics is negative
the other is at least three, so in all cases \(e=e_V+e_{V'}\ge2T\). In
case (2) with \(a=0\), \(b_2(\cD V)=2k>0\) and \(\chi(V)=1+k\), so
\(e\ge e_V\ge\tfrac23T(1+k)\); the case \(a=k\) is symmetric. Both bounds
are at least \(2T\) when \(k\ge2\).
\end{proof}

In an entering class \(k\) counts the \(S^1\times S^2\) factors in
the PSC prime decomposition. Spherical factors can also occur. For
\(k=0\) both sides are rational homology balls. For \(k=1\) their
rational homologies are those of \(S^2\) and \(S^1\).
The winding-zero-first cut of a normalized mixed presentation has
\(k=2\), so it cannot be such a positive-double seam. It can still
receive a nonzero-degree map in Corollary~\ref{cor:marked}.

The sign assumption has a test on each entire side \(U\). With
\(H_{\rm out}\) the trace mean curvature its reflected Yamabe sign
equals the sign of the lowest eigenvalue of
\[
 L=-6\Delta+R,\qquad 3\partial_\nu u+H_{\rm out}u=0.
\]
In a totally geodesic gauge the reflected energy form is the sum of the
two Neumann forms on functions with equal traces. Even reflection attains
the same first eigenvalue, and conformal covariance transfers the sign
back. In any smooth gauge the side form has boundary term
\(2\int_{\partial U}H_{\rm out}u^2\). Without a sign hypothesis one
still has \(e_U\ge T\chi(U)-Y(\cD U)^2/96\) by closed Gauss--Bonnet on
Yamabe approximants \cite{Baker2026ChiralMaster}.
For a positive double one further has
\(e_U\ge T(\chi(U)-1/d)\), with \(d=|\pi_1(\cD U)|\) and \(1/\infty=0\).
For finite \(d\) use the Yamabe upper bound on the universal cover in
Gauss--Bonnet. For infinite \(d\), a strict violation would give a
smooth approximant with positive integrated \(\sigma_2\) and positive
Yamabe sign, hence positive Ricci curvature and finite fundamental group
\cite{ChangGurskyYang2002PositiveRicci,GurskyViaclovsky2003FullyNonlinear}.

Integral homology imposes a separate restriction: the two kernels of the
maps on torsion \(H_1(\Sigma)\) are complementary metabolizers of the
linking pairing. For a whole spherical prime embedded smoothly the
remaining candidates are \(S^3\) and \(S^3/Q_8\); the Poincar\'e sphere
in the locally flat list is excluded by the smooth homology-ball
obstruction \cite[Theorem 5.2]{Hillman2024LocallyFlat}
\cite[Theorem 1.2 and \S8.1]{OzsvathSzabo2003Graded}.
A nontrivial integral homology-sphere PSC seam with positive doubles
requires \(119T/60\le e<2T\) in the entering band. Briefly, it is a
sum of signed Poincar\'e spheres. A finite double group equals the side
group by the normal form for the double amalgam. It is generated by
images of the boundary factors. Triviality would give a contractible
filling, excluded by \cite[Proposition 4.2]{ChodoshMaximoMukherjee2024Complete}.
A nontrivial factor image has order \(60\) or \(120\). Order \(60\) for
the whole group would give \(A_5\), impossible since an integral homology
ball has \(H_2(\pi_1U;\mathbb Z)=0\) by the Hopf sequence, whereas
the nonsplit binary icosahedral central extension and perfection give
\(H_2(A_5;\mathbb Z)\ne0\). Each side therefore costs at least
\(119T/120\). The marked realizations are further constrained by the
boundary isometry and both normal signs
\cite[SEAM, finite covers and paired fillings]{Baker2026MarkedCutCalculations}.

\section{One tube}\label{sec:tube}

\subsection{The form}
Suppose the seam is the pure \(S^1\times S^2\) and its side with
\(b_2=1\) is identified as a normal sphere disk bundle. The trivial
normal Euler number then identifies that side as \(S^2\times D^2\).
Its zero section is a two-knot with exterior \(X\), and \(M\) results
from surgery on the circle in \(X\cup(S^1\times D^3)\). The two
framings differ by the Gluck twist. This describes the specified split
zero-framed lower unknot models of \cite{Baker2026HandleCuts}.

\subsection{Circle-symmetric tubes}
Let \(\Gamma\subset SO(4)\) act freely on \(S^3\), let \(\varphi\) be an
orientation-preserving isometry of \(S^3/\Gamma\) of finite order \(\tau\) with a fixed point
\(p\), and let
\[
 N=\bigl(\mathbb R\times S^3/\Gamma\bigr)\big/\bigl((t+\ell,y)\sim(t,\varphi(y))\bigr)
\]
with the product metric, which is conformally flat. Let \(C\) be the
circle through \(p\), let \(r\) be the distance to \(p\) in \(S^3/\Gamma\),
and let \(\theta\) be smaller than the injectivity radius at \(p\). The
differential of \(\varphi\) at \(p\) is a rotation of angle \(\alpha\),
\(\tau\alpha\in2\pi\mathbb Z\), about an axis. Dividing the metric by
\(\sin^2r\) gives on \(\{r\ge\theta\}\), near the tube, the twisted product
\begin{equation}\label{eq:exterior}
 a_{\rm ext}^2\,d\psi^2+ds^2+d\vartheta^2+\sin^2\vartheta\,(d\phi-\beta\,d\psi)^2,
 \qquad a_{\rm ext}=\frac{\ell}{2\pi\sin r},\quad ds=\frac{dr}{\sin r},
\end{equation}
where \(\psi=2\pi t/\ell\), \((\vartheta,\phi)\) are polar coordinates on
the unit normal sphere about the axis, and \(\beta=\alpha/2\pi+n\) for an
integer \(n\) that depends on the choice of the coordinate \(\phi\). The
two-dimensional factor has curvature \(-1\). We close the tube by a cap
\(S^2\times D^2\) with the metric
\begin{equation}\label{eq:cap}
 g=ds^2+a(s)^2d\psi^2+d\vartheta^2+\sin^2\vartheta\,(d\phi+b(s)\,d\psi)^2,
 \qquad 0\le s\le S,
\end{equation}
and write \(K=-a''/a\) and \(F=b'/a\). Impose smooth polar data at the disk centre, for example
\(a(s)=sP(s^2)\), \(b(s)=s^2Q(s^2)\) with \(P,Q\) smooth and \(P(0)=1\),
and take \(a>0\) on \((0,S]\).

\begin{proposition}[Energy of a cap]\label{prop:cap}
For the metric \eqref{eq:cap},
\(R=2(K+1)-\tfrac12F^2\sin^2\vartheta\) and
\begin{equation}\label{eq:cap-energy}
 e_{\rm cap}=\tfrac12\int|W|^2_{\rm HS}\,dV
 =\pi^2\int_0^Sa\Bigl[\tfrac43\bigl(K+1-\tfrac23F^2\bigr)^2+4F^2
 +\tfrac{16}{135}F^4+\tfrac23F'^2\Bigr]ds .
\end{equation}
The cap glues to the exterior \eqref{eq:exterior} at radius \(\theta\) to a
\(C^{1,1}\) metric exactly when
\[
 a(S)=\frac{\ell}{2\pi\sin\theta},\qquad a'(S)=-\frac{\ell}{2\pi}\cot\theta,
 \qquad F(S)=0,\qquad\int_0^SFa\,ds=-\beta .
\]
The first-jet-matched metric has \(e=e_{\rm cap}\); the smooth
approximating energies tend to this value. Positivity requires control
of the closed conformal Laplacian.
\end{proposition}
\begin{proof}
The formulas for \(R\) and for the integral of \(|W|^2\) over the fibre
are a direct computation with the metric \eqref{eq:cap}. They were
verified symbolically for arbitrary \(a\) and \(b\), and tested against
the Gauss--Bonnet theorem on two closed bundles with varying \(F\); the
script is recorded with the note cited in the manuscript's source record.
The matching conditions are the continuity of the metric and of its first
derivatives; the connection of the exterior is flat, which gives
\(F(S)=0\), and the prescribed signed trivialization gives \(b(S)=-\beta\).
The original exterior is conformally flat. Its rescaling by \(\csc r\)
is prescribed only near the tube, where the rescaled metric has scalar
curvature zero. Extend the positive factor smoothly over the remaining
exterior. Its Weyl curvature is still zero. The extension need not have
nonnegative scalar curvature; a positive rescaling of the cap matching
the original positive exterior, or a direct test of the closed form,
can establish positivity.
\end{proof}

For \(F=0\) the functional \eqref{eq:cap-energy} is
\(\tfrac43\pi^2\int a(1+K)^2ds\): the cap is the product of the unit
sphere with a surface of revolution, which is conformally flat exactly
where \(K=-1\).

\begin{proposition}\label{prop:tube-bounds}
\begin{enumerate}
\item For \(\theta=\pi/2\) the seam is totally geodesic, and every cap
with positive double has \(e_{\rm cap}\ge\tfrac43T\). Equality holds for
\(\ell=2\pi\), \(\beta=0\), \(a=\sin s\): the cap is half of the product
of two unit spheres.
\item The manifold obtained by closing the tube with a cap
\eqref{eq:cap} carries an effective circle action, which on the cap rotates
the disc with weight \(\tau\) and the sphere with weight \(\tau\beta\). If it
is a homotopy four-sphere it is diffeomorphic to \(S^4\).
\end{enumerate}
\end{proposition}
\begin{proof}
(1) The cap has \(\chi=2\) and its double is an \(S^2\)-bundle over
\(S^2\) with \(b_2=2\), so \eqref{eq:doubling} applies; the product of two
unit spheres has \(E_\pm=\tfrac83T\). (2) The flow \(t\mapsto t+s\) on the
exterior has period \(\tau\ell\). In the coordinates of \eqref{eq:exterior}
it is \((\psi,\phi)\mapsto(\psi+2\pi s/\ell,\phi+2\pi\beta s/\ell)\), which
preserves \eqref{eq:cap} and extends over \(s=0\). Over one period the
disc turns \(\tau\) times and the sphere \(\tau\beta=\tau\alpha/2\pi+n\tau\)
times, and these integers are coprime when \(\varphi\) has order exactly
\(\tau\) at \(p\). A homotopy four-sphere with an effective locally smooth
circle action is standard by the work of Fintushel and Pao
\cite{Fintushel1976Circle,Pao1978Nonlinear} together with the Poincar\'e
conjecture for the orbit space \cite{Perelman2003Surgery}.
\end{proof}

When \(S^3/\Gamma\) is the \(\tau\)-fold cyclic branched cover of a knot
and \(\varphi\) generates the covering group, \(N\setminus\) tube is the
exterior of the \(\tau\)-twist-spun knot by Zeeman's theorem
\cite{Zeeman1965Twisting}, \(\alpha=2\pi/\tau\), and the two framings give
the branched twist spins with weights \((\tau,1)\) and \((\tau,1-\tau)\)
in the sense of \cite{Fukuda2018Gluck}.

\begin{remark}[Supporting computations]\label{rem:numerics}
Direct minimization of \eqref{eq:cap-energy} over piecewise profiles gives
admissible caps with \(R>0\) and the following values of \(e/\pi^2\); the
wall \(2T\) is \(8\). Untwisted caps on \(S^1\times S^3\):
\(5.3333\) at \(\theta=\pi/2\), then \(7.4196\), \(7.8475\), \(7.9526\) at
\(\theta=\pi/3,\pi/4,\pi/5\). The flat exterior
\(S^1\times(\mathbb R^3\setminus B)\): above \(8\) at every radius tried.
Two-twist-spun trefoil: \(7.6271\). Two-twist-spun figure-eight knot:
\(7.8025\). Three-twist-spun trefoil: \(7.8928\), and \(7.9736\) for the
other framing. Four-twist-spun trefoil: \(7.9875\) and \(8.1098\).
Five-twist-spun trefoil: \(8.0043\). These values are floating
evaluations of chosen profiles. Each row needs its own feasibility and
closed-sign check before it certifies entry.
\end{remark}

For a different polynomial cap on the two-twist-spun trefoil, take
\(\theta=8\pi/25,\ \ell=16/5,\ \beta=1/2\). Rational interval integration
on 2,048 cells gives the following bounds on the matched metric:
\begin{equation}\label{eq:trefoil}
 7.6822\ \le\ e_{\rm cap}/\pi^2\ \le\ 7.7124<8.
\end{equation}
The positive conformal factor \(w\) on the cap has \(w\ge2259/5000\) and
\((-6\Delta+R)w\ge3951/10000\); it matches the regular positive exterior
to first order. The closed form is strictly positive and smooth
approximants retain the strict energy bound. The join is nonumbilic.
The core is knotted (its exterior group is
\(C_3\rtimes_{-1}\mathbb Z\)); the smooth sphere is standard
\cite{Zeeman1965Twisting,Gordon1976Knots,Baker2026MarkedCutCalculations}.
Within the ansatz the lower bound
\(e_{\rm cap}\ge(16\pi^2/3)(1-a'(S))\) has the matched round equality
profiles. For zero twist and \(a'(S)\le0\), the change of variable
\(u=\int_0^s a\) reduces the fixed-area minimum to a clamped cubic.
Optimization over its terminal area gives the continuum minimum and
excludes every finite untwisted flat-exterior ratio
\cite[TUBE, exact area minimum]{Baker2026MarkedCutCalculations}.

\paragraph{What the exterior recognizes.}
Suppose \(X\) is a smooth compact rational-homology-circle exterior
with connected umbilic marked \(S^1\times S^2\) and positive reflected
class. At \(e_X=0\), reflect the conformally flat atlas in the umbilic
charts. The smooth conformal double has positive sign, \(b_1=1\), and
\(\chi=0\). Bour--Carron's equality theorem identifies its conformal
cylinder cover \cite[Corollary 1.3(iii)]{BourCarron2015Optimal}.
The lifted seams separate the two punctures as equators times a line
in the cylinder model, as shown by the reflection and deck action
\cite[TUBE, zero-energy exterior]{Baker2026MarkedCutCalculations}.
There is a single lifted seam and the side is \(S^1\times D^3\).
Both normal framings give standard \(S^4\)
\cite[Proposition 2]{DingGeiges2010Diffeotopy}.
At positive energy, punctured odd-lens mapping tori can collapse with
\(e_X\downarrow0\) and noncyclic group. In the pinching argument
noncyclicity forces \(Y(\cD X)^2\le96e_X\);
a stronger relative pinching assumption constrains the double and
meridian \cite[Theorems 1.4 and 1.6]{ChenZhu2014Conformal}.
Absolute positive-energy recognition of arbitrary marked homotopy-sphere
fillings remains open.

For the oriented mapping torus of a punctured closed three-manifold,
existence of a positive reflective double is equivalent to the capped
fibre having no aspherical prime. Necessity uses stable representatives
of the fibre class and projection in the cyclic cover. Sufficiency
uses a slow path of positive doubling metrics on that punctured fibre
\cite{Kotschick2017Formal,CarlottoLi2024Boundary,Baker2026MarkedCutCalculations}.
This is an existence statement over metrics; testing \(e_X<2T/3\)
and the sign requires the same metric and the same finite base length.
For the marked \(Q_8\times C_{17}\) prism the ordinary-core slope
\(a\) has orbit orbifold group \(S_3\). The reverse-surgery circle is
isotopic to the surviving order-two orbit, so the even initial framing
has a smooth effective circle action and gives standard \(S^4\).
The odd initial framing gives a homotopy sphere. The odd native slopes
have infinite orbit group and the tested orthogonal cable also fails
the group criterion. A later marked reduction identifies the odd model
as surgery on a section of an \(L(7,2)\#L(7,2)\)-bundle over the circle.
In the fourfold cover the fibre coordinate of the lifted section is a
nontrivial commutator between the two summands, so localization to a
single prime does not give this surgery. The odd smooth type and a subcritical prism
metric remain open
\cite[TUBE, marked surgery and its independent replay]{Baker2026MarkedCutCalculations}.

\paragraph{The projective-plane disk.}
Fix the right-invariant lift on the left quotient \(Q_8\backslash S^3\) with
lengths \((1,1,\sqrt3)\) and collapsing short \(i\)-axis. The boundary
has \(|\eta|=5/12\); an umbilic filling on the normal disk bundle has
\(|E_+-E_-|=5\pi^2\). Hence its half-sum is below \(T\) exactly when
the smaller block costs less than \(3\pi^2/2\). Whole-side Robin-positive
fillings at this fixed marking satisfy the bounds
\[
       \tfrac83\ \le\ \inf e_{\rm disk}/\pi^2\ \le\ 6.46666 .
\]
The upper bound is a certified triaxial trial whose final eighth alone
costs at least \(1.87399\pi^2\) in the smaller block. It does not
decide the strict threshold. A conic orbit-shape family costs at least
\(112\pi^2/27\) in the smaller block, and retaining the specified
positive-scalar orbifold cap's annulus down to \(r=2/3\) already costs
\(3500\pi^2/2187>3\pi^2/2\) there. Switching the two short axes in the
gluing of two standard disk bundles gives the standard smooth sphere.
For the certified improved switched trial the two smooth reflected
signs are positive while the closed energy per chirality lies between
\(12.8999\pi^2\) and \(12.9334\pi^2\).
The long-axis biaxial bound concerns a different disk marking
\cite[SEAM and TUBE, short-axis fillings]{Baker2026MarkedCutCalculations}.

\section{Sweepouts}\label{sec:sweepouts}

\subsection{Level sets near a critical point}
For a symmetric form \(Q\) on \(\mathbb R^4\) with eigenvalues
\(q_1,\dots,q_4\) and unit eigenvectors \(e_i\), let \(s_i(Q)\) be the
second elementary symmetric function of the three eigenvalues other than
\(q_i\).

\begin{lemma}\label{lem:restriction}
For every unit vector \(\nu=\sum\nu_ie_i\),
\(\ \sigma_2\bigl(Q|_{\nu^\perp}\bigr)=\sum_i\nu_i^2\,s_i(Q)\).
\end{lemma}
\begin{proof}
\(\sigma_2(Q|_{\nu^\perp})=\tfrac12\bigl[(\tr Q-Q(\nu,\nu))^2-(|Q|^2-2|Q\nu|^2+Q(\nu,\nu)^2)\bigr]
=\sigma_2(Q)-\tr Q\,Q(\nu,\nu)+|Q\nu|^2\), and
\(\sigma_2(Q)-q_i(\tr Q-q_i)=s_i(Q)\).
\end{proof}

\begin{theorem}[No index two]\label{thm:index}
Let \(f\) be a Morse function on a Riemannian four-manifold, \(p\) a
critical point, \(Q=\Hess f(p)\). For a regular point \(x\) let
\(R_\Sigma(x)\) be the scalar curvature at \(x\) of the level set through
\(x\) with its induced metric.
\begin{enumerate}
\item If \(s_i(Q)<0\) for some \(i\), then \(R_\Sigma\) is unbounded below
near \(p\).
\item If \(s_i(Q)>0\) for all \(i\), then \(R_\Sigma(x)\to+\infty\) as
\(x\to p\).
\item The inequalities \(s_i(Q)\ge0\), \(i=1,\dots,4\), hold at every
critical point of index \(0\) or \(4\), fail at every critical point of
index \(2\), and hold at a critical point of index \(1\) or \(3\) exactly
when the eigenvalue \(q\) of exceptional sign satisfies
\(|q|^{-1}\ge|q_j|^{-1}+|q_k|^{-1}\) for every pair of the other three.
\end{enumerate}
\end{theorem}
\begin{proof}
In normal coordinates at \(p\), \(\nabla f(x)=Qx+O(|x|^2)\) and
\(\Hess f(x)=Q+O(|x|)\). The level set through \(x\) has unit normal
\(\nu=Qx/|Qx|+O(|x|)\) and second fundamental form
\(\pm\Hess f|_{\nu^\perp}/|\nabla f|\). By the Gauss equation and
Lemma~\ref{lem:restriction},
\[
 R_\Sigma(x)=R-2\Ric(\nu,\nu)+\frac{2}{|Qx|^2}\Bigl(\sum_i\nu_i^2s_i(Q)+O(|x|)\Bigr).
\]
Along \(x=te_i\) the normal tends to \(\pm e_i\) and the bracket to
\(s_i(Q)\); this gives (1). If all \(s_i\) are positive the bracket is
bounded below by a positive constant near \(p\), which gives (2). For (3):
if all eigenvalues have one sign every \(s_i\) is positive. Let the index
be two, with eigenvalues \(\alpha_1,\alpha_2>0>-\beta_1,-\beta_2\). The
four conditions are \(\beta_1\beta_2\ge\alpha_i(\beta_1+\beta_2)\) and
\(\alpha_1\alpha_2\ge\beta_j(\alpha_1+\alpha_2)\). The first two give
\(\alpha_i<\min\beta\) and the last two \(\beta_j<\min\alpha\), which is
absurd. For index one with eigenvalues \(\alpha_1,\alpha_2,\alpha_3>0>-\beta\),
omitting \(-\beta\) gives a positive value and omitting \(\alpha_i\) gives
\(\alpha_j\alpha_k-\beta(\alpha_j+\alpha_k)\). Index three is the same
after changing the sign of \(f\).
\end{proof}

\begin{corollary}[Recognition by a sweepout]\label{cor:sweepout}
Let \(M\) be a closed connected oriented four-manifold with a metric and a
Morse function for which one finite lower bound holds for the scalar curvature
at every point of every regular level. Then \(M\) is diffeomorphic to \(S^4\) or to a connected sum of
copies of \(S^1\times S^3\). For a smooth homotopy four-sphere \(M\) the
following are equivalent: \(M\cong S^4\); for some metric there is a Morse
function whose regular levels have scalar curvature bounded below; for
some metric there is a Morse function whose regular levels have positive
scalar curvature.
\end{corollary}
\begin{proof}
By Theorem~\ref{thm:index} the function has no critical point of index
two. Rearranging, \(M\) is the union of the handles of index at most one
and the handles of index at least three, two handlebodies
\(\natural_kS^1\times B^3\) glued along their boundary. By the theorem of
Laudenbach and Po\'enaru \cite{LaudenbachPoenaru1972Handlebodies}, and by
Cerf's theorem \cite{Cerf1968Gamma4} for \(k=0\), every diffeomorphism of
the boundary extends over one side, so \(M\) is the double. The height
function of the round sphere gives the converse.
\end{proof}

Strictness in Theorem~\ref{thm:index} matters. In Euclidean
coordinates \((x,y,z,w)\), set
\(f=x^2+y^2+z^2-w^2/2+cxw^2\). On \((x,0,0,0)\) its regular scalar
curvature is \(4c/x\); the pure quadratic with the same Hessian has
nonnegative regular scalar curvature. In an entering gauge
\(R_\Sigma=P+2\sigma_2(B)\), so \(\sigma_2(B)\ge0\) suffices. For every
smooth metric on an identified standard sphere a two-extrema height has
a uniform finite lower bound: its Hessians are definite at the two
poles and compactness controls the remaining band. Strict positivity
on every regular level in a specified entering gauge is a stronger
condition. The topology of the levels alone has the following limitation.

\begin{proposition}[Presentations]\label{prop:presentation}
Let \(\langle x_1,\dots,x_\ell\mid r_1,\dots,r_\ell\rangle\) be a
presentation of the trivial group and let \(M_P\) be obtained from
\(\#_\ell S^1\times S^3\) by surgery on disjoint embedded circles in the
classes \(r_1,\dots,r_\ell\), with either oriented normal framing
on each circle. Then \(M_P\) is a smooth homotopy
four-sphere, and it has a Morse function whose regular level sets are all
connected and diffeomorphic to connected sums of copies of
\(S^1\times S^2\).
\end{proposition}
\begin{proof}
Write \(\#_\ell S^1\times S^3=N_0\cup N_1\), two copies of
\(\natural_\ell S^1\times B^3\), and isotope the circles into
\(\partial N_0\). Surgery on a circle \(K_i\) adds to the handles of the
manifold a two-handle along \(K_i\) and a two-handle along a zero-framed
meridian \(U_i\) of \(K_i\). So \(M_P=N_0\cup\{U_i,K_i\}\cup N_1\), with
fundamental group \(\langle x\mid r\rangle=1\) and Euler characteristic
two. Attach the handles in the order \(U_1,\dots,U_\ell,K_1,\dots,K_\ell\).
In the
boundary surgery picture take each \(U_i\) a small zero-framed meridian.
Its disk meets only \(K_i\), once. After zero surgery a parallel
punctured disk capped by a meridian disk of the attached solid torus
gives a disjoint sphere meeting \(K_i\) once. The level after all
\(U_i\) is \(\#_{2\ell}(S^1\times S^2)\). For the next \(K_i\), its
dual sphere punctured along the knot gives a compressing disk for the
boundary of the knot exterior with meridional boundary. A neighborhood
of this disk and the boundary torus is a punctured solid torus.
Separating it off by its spherical boundary, meridional filling adds an
\(S^1\times S^2\) summand, while any integer framing intersects the
meridian once and adds an \(S^3\) summand. Thus each \(K_i\)-surgery
removes one \(S^1\times S^2\) factor; the other dual spheres survive.
The levels within the two outer handlebodies have the asserted form.
\end{proof}

When the presentation is trivial under the moves of Andrews and Curtis,
\(M_P\) is standard; in general this is not known
\cite{GompfScharlemannThompson2010}. So the diffeomorphism types of all
levels of a Morse function, even in a prescribed order, do not recognize
the sphere with present knowledge, while the geometry of the levels near
the critical points of index two does, by Theorem~\ref{thm:index}. The
recognition theorem of \cite{Baker2026HandleCuts} concerns presentations
with two two-handles, where one intermediate level remains.

\subsection{The minimal cut}
An entering gauge has positive Ricci curvature. By Zhou's theorem
\cite{Zhou2015MinMax} the min-max hypersurface of the fundamental class is
orientable with multiplicity one and index one in our case.
Indeed every closed embedded hypersurface in a homotopy sphere is
two-sided by mod-two separation. This excludes the alternative with
multiplicity two on a nonorientable support; stability there refers to
one-sided normal variations \cite[Theorem 1.1]{Zhou2017MinMax}.

\begin{proposition}\label{prop:index-one}
Let \(\Sigma\) be a closed two-sided minimal hypersurface of Morse index
at most one in a Riemannian four-manifold with \(R>0\). Then the conformal
Laplacian \(-8\Delta+R_\Sigma\) of \(\Sigma\) has at most one nonpositive
eigenvalue.
\end{proposition}
\begin{proof}
The Jacobi operator is \(J=-\Delta-(\Ric(\nu,\nu)+|B|^2)\), and the Gauss
equation with \(H=0\) gives \(\Ric(\nu,\nu)+|B|^2=\tfrac12(R-R_\Sigma+|B|^2)\).
Hence, as quadratic forms,
\(-8\Delta+R_\Sigma\ge2J+R+|B|^2\). The right side is positive on a
subspace of codimension at most one.
\end{proof}

In a smooth neighborhood of the unit round metric on the identified
standard \(S^4\), every integral fundamental-class first-width
hypersurface is the zero level of a two-extrema Morse function with
strictly positive scalar curvature on every regular level.
To see this, take any sequence approaching round. Round latitudes
bound the widths from above. Bounded area and index give single-sheet
smooth graphical convergence after a subsequence
\cite[Appendix A]{Sharp2017Compactness}. The round limit is an
equator. On a round minimal hypersurface some normal coordinate
\(\langle v,\nu\rangle\) of a constant Euclidean vector is a nonzero
Jacobi eigenfunction of eigenvalue \(-3\) by Codazzi and minimality.
The constant Rayleigh test is strictly below \(-3\) unless \(B=0\);
otherwise the index is at least two. See
\cite[GEO, first-width alignment]{Baker2026MarkedCutCalculations}.
Choose in the five-dimensional eigenspace cluster near four a unit
height minimizing its squared restriction to the given hypersurface.
The limiting restriction form has precisely the normal height as its
kernel. The cut is a small graph over that height's zero set. Subtract
the graph height with a cutoff in a fixed band. The resulting height
is \(C^2\)-close to the round height; its regular leaves have positive
scalar curvature in the middle by compactness and at the ends by the
definite Hessian and \(P>0\). No uniqueness of the width surface is used.

A contrasting local test uses the Berger orbit with squared fibre
length \(65/16\) and \(B=\mathrm{diag}(1/\sqrt8,-1/\sqrt8,0)\) in its
orthonormal frame. Its intrinsic scalar is \(-1/8\), its Jacobi index
is one, and its ambient germ has \(R>0,\ \sigma_2(A)>0\).
The marked eta and extrinsic contribution give side differences
\(-75\pi^2,+75\pi^2\) on a closed homotopy sphere. Equality in the
resulting lower bound would force roundness by Section~\ref{sec:ledger},
so \(e>75\pi^2\). Consequently this germ has no closed entering
completion. For the full specified germ, a short-axis interchange
gives a smooth twisted double of each cap. Gauss--Bonnet on that
double gives \(4\int_{V_i}\sigma_2(A)<-67\pi^2\) on each cap. In
a specific completed test metric an angular sweepout has
maximal area at most \(1.903738\), strictly below the central orbit area
\(\pi^2\sqrt{65}/2\); this metric has a negative-scalar orbit on each
cap \cite[GEO, Berger completion]{Baker2026MarkedCutCalculations}.
A local negative example therefore leaves the induced first eigenvalue
of an unrestricted first-width hypersurface in a closed entering gauge
undecided.

\section{Free cuts}\label{sec:free}

\begin{proposition}[Descent along a free cut]\label{prop:free}
Let a free cut split a smooth homotopy four-sphere as
\(M=V_1\cup_\Sigma V_2\). Close each side by a round hemisphere with
matching first jet, writing \(\widehat V_i\) for the resulting closed
smooth manifold. Then:
\begin{enumerate}
\item Each \(\widehat V_i\) is a homotopy sphere,
\(M\cong\widehat V_1\#\widehat V_2\), and \(e=e_{V_1}+e_{V_2}\).
Both chiral energies of the matched metric on \(\widehat V_i\) equal
\(e_{V_i}\). There are smooth metrics with energies tending to this value.
\item If both matched capped classes have strictly positive Yamabe sign
and \(e<2T\), the lower-energy capped summand is standard. If both
\(e_{V_i}<T\), then \(M\cong S^4\).
\item If in addition the original smooth positive class attains the
infimum over positive classes on its smooth type, then the minimizing class is conformally round and \(e=0\).
\end{enumerate}
\end{proposition}
\begin{proof}
Normalize the boundary value and normal derivative of a conformal
factor so the induced metric is round and \(B=0\). The boundary
orientation can be chosen to match either hemisphere. The cut has
zero eta and zero extrinsic term, so the chiral integrals on each
side agree. A hemisphere adds zero Weyl energy. Van Kampen and
Mayer--Vietoris across \(S^3\) make both sides simply connected
homology balls; capping gives smooth homotopy spheres. The boundary
gluing extends over a standard ball \cite{Cerf1968Gamma4}.

Smoothing a piecewise smooth \(C^{1,1}\) matched metric in collar
charts yields convergence in \(C^1\cap W^{2,p}\) for finite \(p\).
Chiral curvature converges in \(L^2\). The closed scalar forms converge
in norm on \(H^1\) since \(H^1\hookrightarrow L^4\) and scalar curvature
converges in \(L^2\). Thus a strictly positive limiting Yamabe sign
persists in the smooth approximants. At least one side has limiting
energy strictly below \(T\); use the smooth positive approximants and
the strict sphere theorem \cite{ChangGurskyYang2003}. This proves (2).
Removing the standard summand \(i\) leaves a
manifold diffeomorphic to \(M\). Its positive approximants give
\(e\le e_{V_j}=e-e_{V_i}\) at attainment, hence \(W=0\) on the open
side \(V_i^\circ\). Positivity is open under smooth metric variations,
and the signature fixes the difference of the two Weyl functionals.
The smooth minimizer is therefore Bach-flat on all of \(M\).
Unique continuation \cite[Theorem 3.8]{LassasLiimatainen2019ConformalHarmonic}
gives \(W=0\) everywhere. Conformal development identifies \((M,[g])\) with the round conformal sphere.
\end{proof}

There is a sufficient cut measurement for the two signs used in
Proposition~\ref{prop:free}. If one smooth representative has \(R>0\)
and the same free cut satisfies
\[
 \max\bigl\{\|H_+\|_{L^3(h)},\|H_-\|_{L^3(h)}\bigr\}
 \le 3(2\pi^2)^{1/3},
\]
where \(H_\pm\) are the positive and negative parts of the trace in one
normal convention, both matched conformally flat caps have strictly
positive sign. The ball trace inequality
\cite[Theorem 1.1]{NgoNguyenPhan2020Trace}
bounds the cap Robin form below by \(6(2\pi^2)^{1/3}\|f\|_3^2\)
for the common boundary value \(f\). This pays the adverse side term
\(2\int H_{\rm out} f^2\) on either side. At equality the positive side
bulk controls any nonzero restriction there, and a cap function with
zero trace has positive Dirichlet energy. The compact form attains its
first eigenvalue. The marked cap construction and conformal transformation
of its outward normal derivative are given in \cite{Baker2026ChiralMaster}.
Thus at attainment the same measurable hypothesis implies conformal roundness.

More generally, let a spherical cut have two first-jet-matched ball
fillings with costs \(c_i=(E_+(\mathrm{cap}_i)+E_-(\mathrm{cap}_i))/2\),
in the actual orientations. The two capped limits now have energies
\(e_{V_i}+c_i\). With both capped signs positive and
\(e+c_1+c_2<2T\) the same argument gives a standard summand \(i\);
it gives both if each individual sum is below \(T\). If the original
type is nonstandard and \(e\le I^+(M)+\delta\), the surviving side
\(j\) is a competitor on \(M\), giving
\[
 e_{V_i}\le c_j+\delta,\qquad
 E_\pm(V_i)\le E_\pm(\mathrm{cap}_j)+\delta.
\]
The explicit cut-off round hemisphere with extension
\(q_i=\chi(s)(h-\gamma\ \pm 2s B)\) in a fixed trace-free boundary
gauge has \(c_i\le3750\pi^2 D_i^2\) if
\(D_i=\max_{k\le2}\sup|\nabla_0^k q_i|\le1/100\).
Here the sign is fixed by the physical normal and \(\chi\) is the
specified cutoff in the hemisphere trial
\cite[SEAM, quantitative capped descent]{Baker2026MarkedCutCalculations}.

\section{Outstanding problems}\label{sec:open}

For the reduced mixed presentations, the first recognition question is
whether the entering geometry supplies a closed positive source
mapping with nonzero degree to one of the two actual intermediate
cuts. Equation~\eqref{eq:marked-average} is one way to certify this
on a native collar. An unrelated positive hypersurface gives no such
map. The final \(S^1\times S^2\) level cannot map with nonzero degree
to the winding-zero cut, whose first Betti number is two, and a degree-one
map to the winding-one homology sphere would already force its perfect
fundamental group to be cyclic and thus trivial. The attaching tubes,
upper closing, and the placement of the source in the band matter.

The sphere statement for all entered types still requires a general
existence or reduction theorem. The first-width construction is available
near the round metric. Away from this regime a positive first-width
hypersurface would still need a connection to the marked target or a
sweepout of both sides. An invariant orbit, a cylindrical singularity,
or a local index computation does not supply either identification.
In the round join the minimal principal orbit has full index six.
The certified trefoil join has only one-sided second jets and a change
of conformal gauge changes its radial minimal leaves. A general
continuation through singularities with the needed regular-level
bounds remains open.

For entry, the atlas and ledger test the proposed pieces in their actual
markings. The raw-energy margin, the closed conformal-Laplacian sign
and the two reflective signs are separate quantities. The calibrated
circle-symmetric trefoil and even prism do not supply a new smooth type.
The positive-energy exterior problem, the strict short-axis cap budget,
and the supply of marked, positive, low-cost spherical caps for descent
remain open in their respective generalities.
\bibliographystyle{unsrturl}
\bibliography{references}
\end{document}
